	
	
\Bsec{Zahlenbasen}{
	
		\begin{itemize}
			\item Umwandlung von dezimalen Zahlen in binäre Zahlen
			\itemf Umwandlung von binären Zahlen in Hex
			\itemf Gleitkommazahlen und ihre Umwandlung
			\itemf IEEE-Standard
			\itemf NSDR Format
		\end{itemize}
		}
		
%----------------------------------------------------------------------------------
\bsubsec{Dezimal zu Binär}{
	Führe folgende dezimale Zahlen ins binäre System um:
		\begin{enumerate}
			\itemf 14
			\itemf 63
			\itemf 1.25
			\itemf $\nicefrac{3}{8}$
			\itemf $\nicefrac{13}{64}$
		\end{enumerate}
%			12 (Dezimal) soll als binäre Zahl geschrieben werden:
%			12 - 2^{3} = 4 \imp 1\\
%			4 - 2^{2} = 0 \imp 1\\
%			0 - 2^{1} < 0 \imp 0\\
%			0 - 2^{0} < 0 \imp 0
%			Somit ergibt sich (von oben nach unten) ''$\fat{1100}_2$''!
		}
%		Führe folgende dezimale Zahlen ins binäre System um:
%		
%		}
	\bsol{
		\begin{enumerate}
			\itemf $1110_2$
			\itemf $111111_2$
			\itemf $01.01_2$
			\itemf $0.011_2$
			\itemf $0.001101_2$
		\end{enumerate}
		}
%----------------------------------------------------------------------------------
\bsubsec{Binär zu Hex}{
		Führe folgende binäre Zahlen ins Hexadezimale System um:
		\begin{enumerate}
			\itemf $10_2$
			\itemf $101010_2$
			\itemf $111000_2$
			\itemf $0.11_2$ (\red{Advanced!})
		\end{enumerate}
	}
%		
%		Führe folgende binäre Zahlen ins Hexadezimale System um:
%		
%	}
	\bsol{
		
		\begin{enumerate}
			\itemf 0x2
			\itemf 0x2A
			\itemf 0x38
			\itemf 0x0.C
		\end{enumerate}
	}
%----------------------------------------------------------------------------------
\bsubsec{Gleitkommazahlen}{
%		
%		Gleitkommazahl: $(-1)^{Vorzeichen} \cdot 2^{Exponent} \cdot Mantisse_2$
%		
%		$2^1 \cdot 10_2 \Longleftrightarrow 4$
%		
%		$2^{-1} \cdot 100_2 \Longleftrightarrow 2$
%	}
%		
		Wie sähen folgende Gleitkommazahlen nach allgemeinem Standard ''normiert'' aus und zu welcher Zahl im gewohnten (dezimalen) System sind sie äquivalent?
		\begin{enumerate}
			\itemf $2^3 \cdot 101_2$
			\itemf $2^{9} \cdot 11_2$
			\itemf $2^{-2} \cdot 101_2$
			\itemf $2^{-4} \cdot 1011.1_2$
		\end{enumerate}
%		
	}
	\bsol{
%		
		\begin{enumerate}
			\itemf $2^5 \cdot 1.01_2 = 101000_2 = 40$
			\itemf $2^{10} \cdot 1.1_2 = 11000000000_2 = 1536$
			\itemf $2^0 \cdot 1.01_2 = 1.01_2 = 1.25$
			\itemf $2^{-1} \cdot 1.0111_2 = 0.10111_2 = \nicefrac{23}{32}$
		\end{enumerate}
	}

%----------------------------------------------------------------------------------
\bsubsec{Maschinenzahlen (IEEE)}{
		
%	}
%		
		Forme folgende Zahlen in eine Maschinenzahl nach IEEE-Standard um:
		\begin{enumerate}
			\itemf $2^3 \cdot 101_2$
			\itemf $-1 \cdot 2^{9} \cdot 11_2$
			\itemf $2^{-2} \cdot 101.11010110001_2$
			\itemf $2^{-4} \cdot 1011.1_2$
		\end{enumerate}
%		
	}
	\bsol{
%		
		\begin{enumerate}
			\itemf $2^5 \cdot 1.01_2 =$\\$0|10000100|01000000000000000000000$
			\itemf $2^{10} \cdot 1.1_2 =$\\$1|10001001|10000000000000000000000$
			\itemf $2^0 \cdot 1.0111010110001_2 =$\\$0|01111111|01110101100010000000000$
			\itemf $2^{-1} \cdot 1.0111_2 =$\\$0|01111110|01110000000000000000000$
		\end{enumerate}
	}
	
%----------------------------------------------------------------------------------
\bsubsec{NSDR (New Super Digit Representation)}{
	
	Eine neue (erfundene) Zahlendarstellung, die ihr für die nächste Aufgabe benötigt
	\begin{itemize}
		\item Besteht aus 10 Bits
		\item Exponent: 5 Bit (verhält sich wie in IEEE)
		\item Mantisse: 5 Bit (Normiert)
		\item 10 Nullen ist vor-reserviert für Infinity
		\item 10 Einser ist vor-reserviert für -Infinity
	\end{itemize}
}

	\orangeSlide{NSDR Tasks}{
		
		Forme folgende Zahlen in eine Maschinenzahl nach NSDR-Standard um:
		\begin{enumerate}
			\itemf $2^3 \cdot 101_2$
			\itemf $-1 \cdot 2^{0} \cdot 11_2$
			\itemf $2^{-2} \cdot 101.1101_2$
			\itemf $2^{11} \cdot 111111.111111_2$
		\end{enumerate}
	}
	\bsol{
		
		\begin{enumerate}
			\itemf $2^5 \cdot 1.01_2 = 10100|10100$
			\itemf Kann nicht dargestellt werden, gibt kein Vorzeichen-Bit o.ä.
			\itemf $2^0 \cdot 1.011101_2 = 01111|10111$ \\(Hier musste gerundet werden, nicht genug Bits)
			\itemf $2^{16} \cdot 1.1111_2 = -\infty$ \\(Nur Einser ist vor-reserviert, ein PC würde vmtl. eine Exception werfen, da er aufrunden müsste und das außerhalb des darzustellbaren Bereiches wäre)
		\end{enumerate}
	}
%----------------------------------------------------------------------------------
